\section{Scripts de arranque}
\label{sec:scripts_arranque}

Para simplificar el despliegue y evitar errores de configuración manual, se desarrollaron
scripts de arranque que unifican en un único comando el lanzamiento de todos los
componentes del sistema. Cada escenario de despliegue tiene su propio script.

\subsection{Escenario monopuesto (\texttt{start\_studio\_single\_pc.sh})}

Este script arranca el sistema completo en un único equipo: inicia voctocore, lanza el
servicio de telemetría, inyecta las señales de continuidad (vídeo de pausa, vídeo
\textit{offline} y música de fondo), y finalmente abre voctogui. Antes de lanzar cada
componente, un bucle comprueba con \texttt{ss -lnt} que el puerto 9999 está en escucha,
evitando las condiciones de carrera que se producían con los \texttt{sleep} fijos del
código original. Al cerrar voctogui, la instrucción \texttt{trap} captura la señal de
salida y termina todos los procesos hijos de forma limpia.

\subsection{Escenario dos PCs (\texttt{2pc\_escenario2/})}

El escenario de dos equipos se gestiona con dos scripts independientes que deben
ejecutarse en paralelo:

\begin{itemize}
  \item \textbf{\texttt{start\_voctocore\_pc1.sh}} (PC servidor): arranca voctocore,
        las señales de continuidad y el servicio de telemetría en el equipo que actúa
        como servidor. Detecta automáticamente la IP de la interfaz LAN mediante la
        tabla de enrutamiento del kernel
        (\texttt{ip route get 1.1.1.1}), de modo que no es necesario configurar la IP
        a mano en cada despliegue.
  \item \textbf{\texttt{start\_voctogui\_pc2.sh}} (PC realizador): lanza voctogui y
        lo apunta a la IP del servidor mediante la variable de entorno
        \texttt{IP\_SERVER}. Si no se especifica, el valor por defecto es
        \texttt{127.0.0.1} para pruebas en un solo equipo.
\end{itemize}

Ejemplo de uso desde el PC realizador:
\begin{verbatim}
IP_SERVER=192.168.1.50 ./2pc_escenario2/start_voctogui_pc2.sh
\end{verbatim}

\subsection{Escenario Docker (\texttt{launch\_docker\_studio.sh})}

Este script ejecuta \texttt{xhost +local:docker} para permitir que los contenedores
accedan al servidor gráfico del anfitrión, exporta la IP local como variable de entorno
(\texttt{HOST\_IP}) para que el servicio de telemetría la incluya en los registros, y
lanza el stack con \texttt{docker compose up}. Al cerrar el terminal, la instrucción
\texttt{trap} ejecuta \texttt{docker compose down} de forma automática.

El mismo flujo está disponible a través del \texttt{Makefile}:
\begin{verbatim}
make docker-up      # xhost + y docker compose up
make docker-down    # detener todos los servicios
make docker-restart # down + up en un solo paso
\end{verbatim}

\subsection{Escenario Kubernetes (\texttt{launch\_k8s.sh})}

El script de Kubernetes comprueba si Minikube y el pod principal ya están en ejecución
antes de lanzarlos, aplica los manifiestos con \texttt{kubectl apply}, espera a que
todos los contenedores estén listos y establece un \texttt{kubectl port-forward} que
expone los puertos del pod en \texttt{localhost} del anfitrión. Esto permite que
voctogui se conecte al clúster exactamente igual que en cualquier otro escenario. Al
salir, un \texttt{trap EXIT} cierra el \textit{port-forward} de forma limpia.

\begin{table}[H]
  \centering
  \caption{Scripts de arranque por escenario de despliegue.}
  \label{tab:scripts_arranque}
  \small
  \begin{tabular}{|l|l|l|}
    \hline
    \textbf{Escenario} & \textbf{Script} & \textbf{Comando de inicio} \\
    \hline
    Un PC (nativo)   & \texttt{start\_studio\_single\_pc.sh}      & \texttt{./start\_studio\_single\_pc.sh} \\
    Dos PCs (PC1)    & \texttt{2pc\_escenario2/start\_voctocore\_pc1.sh} & \texttt{./start\_voctocore\_pc1.sh} \\
    Dos PCs (PC2)    & \texttt{2pc\_escenario2/start\_voctogui\_pc2.sh}  & \texttt{IP\_SERVER=<IP> ./start\_voctogui\_pc2.sh} \\
    Docker Compose   & \texttt{launch\_docker\_studio.sh}         & \texttt{make docker-up} \\
    Kubernetes       & \texttt{launch\_k8s.sh}                    & \texttt{./launch\_k8s.sh} \\
    \hline
  \end{tabular}
\end{table}
